
\section*{Part II. Applications of optics \hfill 8 / 15}
\subsection*{1. Postselection changes the quantum operation}
If a detector outcome \(k\) is retained, the unnormalised
state and probability are
\[
\tilde\rho_k=K_k\rho K_k^\dagger,\qquad
p_k=\operatorname{Tr}\tilde\rho_k,\qquad
\boxed{\rho_k=\tilde\rho_k/p_k.}
\]
In contrast, \(\sum_k\tilde\rho_k\) is trace-preserving.
For a weakly reflecting beam splitter with a vacuum ancilla,
detecting \emph{one} reflected photon implements
conditional photon subtraction. For input \(\ket n\),
\[
K_{1}\ket n=
\sqrt{n(1-\eta)}\,\eta^{(n-1)/2}\ket{n-1}
\]
up to an irrelevant overall phase. Higher-\(n\) amplitudes
are preferentially weighted by \(\sqrt n\); this is
\emph{not} the same as deterministically deleting a photon.
A no-click conditioning has different photon-number weights.
\begin{methodbox}
\textbf{Exam method:} transform the joint input with the beam
splitter; project the detector mode onto the specified Fock
state; retain the unnormalised signal; calculate \(p\);
finally normalise.
\end{methodbox}

\subsection*{2. Jaynes--Cummings: one atom plus one cavity mode}
Near resonance, the rotating-wave approximation gives
\[
\boxed{H_{\rm JC}=\hbar\omega_c a^\dagger a+
\hbar\omega_a\sigma_+\sigma_-+
\hbar g(a\sigma_++a^\dagger\sigma_-).}
\]
The total excitation number is conserved, so
\(\{\ket{e,n},\ket{g,n+1}\}\) form independent
two-dimensional subspaces. At resonance their
symmetric/antisymmetric dressed-state energies are split by
\(2\hbar g\sqrt{n+1}\). In the one-excitation manifold the
vacuum Rabi doublet has frequency separation \(2g\).
Away from resonance, diagonalise
\(\frac{\hbar}{2}\begin{pmatrix}\Delta&2g\sqrt{n+1}\\
2g\sqrt{n+1}&-\Delta\end{pmatrix}\)
after subtracting a common energy.

\subsection*{3. Tavis--Cummings and bright/dark states}
For \(N\) identical atoms coupling in phase to a common mode,
\[
H_I=\hbar g(aJ_++a^\dagger J_-),\quad
J_\pm=\sum_i\sigma_\pm^{(i)}.
\]
For two atoms, the symmetric one-excitation state
\(\ket S=(\ket{eg}+\ket{ge})/\sqrt2\) couples to
\(\ket{gg,1}\) with strength \(\sqrt2g\);
the antisymmetric \(\ket A=(\ket{eg}-\ket{ge})/\sqrt2\)
is dark under identical coupling. For a symmetric single
excitation in \(N\) atoms, the coupling is \(g\sqrt N\).
This collective enhancement underlies superradiant behaviour,
though collective coherent exchange and spontaneous superradiant
decay are not identical phenomena.

\src{W5L3 pp.~3--25; W6L1 pp.~5--29; W6L2 pp.~5--25.}
